%!LW recipe=latexmk (lualatex)

\documentclass[a4paper,12pt]{article}

% =========================
% układ strony
% =========================
\usepackage[
  paper=a3paper,
  left=1.35cm,
  right=1.35cm,
  top=1.2cm,
  bottom=1.55cm,
  footskip=18pt
]{geometry}

\pagestyle{plain}
\raggedbottom

% =========================
% typografia i język
% =========================
\usepackage{microtype}
\usepackage{fontspec}
\usepackage{polyglossia}
\setdefaultlanguage{polish}
\setmainfont{STIX Two Text}

% =========================
% matematyka
% =========================
\usepackage{amsmath,amssymb,mathtools}
\usepackage{unicode-math}
\usepackage{diagbox}
\setmathfont{STIX Two Math}

% =========================
% dodatki
% =========================
\usepackage{xcolor}
\usepackage{multicol}
\usepackage{enumitem}
\usepackage[most]{tcolorbox}
\usepackage{booktabs} % tabelki
\usepackage{tikz}
\usetikzlibrary{automata, positioning, arrows.meta}
% =========================
% kolory
% =========================
\definecolor{TaskGray}{HTML}{F7F7F8}

% =========================
% globalna zwartość / odstępy
% =========================
\setlength{\parindent}{0pt}
\setlength{\parskip}{1.6pt}
\setlength{\columnsep}{0.9cm}
\setlength{\columnseprule}{0.5pt}
\linespread{1.02}

% =========================
% itemize, steps i podpunkty 
% =========================

\setlist[itemize]{
  label=\textbullet,
  leftmargin=1.2em,
  itemsep=0pt,
  topsep=0pt,
  parsep=0pt,
  partopsep=0pt
}

\newlist{steps}{enumerate}{1}
\setlist[steps]{
  label=\textbf{(\arabic*)},
  leftmargin=1.9em,
  itemsep=0.42em,
  topsep=0.2em,
  parsep=0pt,
  partopsep=0pt
}

\newlist{podpunkty}{enumerate}{1}
\setlist[podpunkty]{
  label=\textbf{(\alph*)},
  leftmargin=1.9em,
  itemsep=0.42em,
  topsep=0.2em,
  parsep=0pt,
  partopsep=0pt
}


% =========================
% indeks
% =========================
\newcommand{\indeks}[1]{%
  \par\bigskip
  \noindent
  {\setlength{\fboxsep}{8pt}%
   \setlength{\fboxrule}{1.3pt}%
   \fcolorbox{black}{white}{%
     \strut\fontsize{22}{26}\selectfont\bfseries #1%
   }}%
  \par\vspace{0.95em}
}

% =========================
% temat i bit
% =========================
\newcommand{\temat}[1]{%
  \par\medskip
  \noindent{\bfseries
      \fontsize{16}{20}\selectfont
      #1%
  }
  \par\medskip
}

\newcommand{\bit}[1]{%
  \par\medskip
  \noindent{\bfseries #1}%
  \par\smallskip
}

% =========================
% logbox
% =========================
\newtcolorbox{logbox}{
  enhanced,
  colback=TaskGray,
  colframe=white,
  boxrule=0.9pt,
  arc=0pt,
  left=6pt,
  right=6pt,
  top=5pt,
  bottom=5pt,
  boxsep=0pt,
  before skip=0.2em,
  after skip=0.7em
}

% =========================
% tw i def
% =========================
\newcommand{\twierdzenie}[2]{%
  \noindent
  #1
  \par\smallskip
  \noindent\hspace*{1.2em}$\Longrightarrow$
  \par\smallskip
  \noindent
  #2
  \par\smallskip
}

\newcommand{\definicja}[2]{\textbf{def #1}$\ \coloneqq\ $#2}
% =========================
% makra matematyczne
% =========================
\newcommand{\N}{\mathbb{N}}
\newcommand{\Z}{\mathbb{Z}}
\newcommand{\Q}{\mathbb{Q}}
\newcommand{\R}{\mathbb{R}}
\newcommand{\C}{\mathbb{C}}
\newcommand{\K}{\mathbb{K}}
\newcommand{\X}{\mathbb{X}}

\DeclareMathOperator{\Lin}{Lin}

\DeclarePairedDelimiter{\abs}{\lvert}{\rvert}
\DeclarePairedDelimiter{\norm}{\lVert}{\rVert}
\DeclarePairedDelimiter{\set}{\lbrace}{\rbrace}
\DeclarePairedDelimiter{\inner}{\langle}{\rangle}

\begin{document}

\bit{skrypt pod kartkówki, wykłady z 25/26}

\bigskip
\indeks{w1-2}

\bigskip
\temat{topologia}

\bigskip
\textbf{def} \; \(X\) - zbiór, \; \(\mathcal{T} \subseteq 2^X\), \; \(\mathcal{T}\) - \textbf{topologia} \, \(\iff\)
\begin{enumerate}[label=\roman*), leftmargin=1.9em, itemsep=0.2em, topsep=0.2em, parsep=0pt, partopsep=0pt]
    \item \(X, \emptyset \in \mathcal{T}\)
    \item \(\forall_{\alpha \in A} \;\, U_\alpha \in \mathcal{T} \implies \displaystyle\bigcup \{U_\alpha : \alpha \in A\} \in \mathcal{T}\)
    \item \(U, \hat{U} \in \mathcal{T} \implies U \cap \hat{U} \in \mathcal{T}\)
\end{enumerate}

\bigskip
\textbf{def} \; \(A\) - \textbf{otwarty} \, \(\iff \, A \in \mathcal{T}\)

\smallskip
\textbf{def} \; \(A\) - \textbf{domknięty} \, \(\iff \, A^c \in \mathcal{T}\)

\bigskip
\begin{logbox}
\textbf{przykłady}
\smallskip
\begin{steps}
    \item \(\mathcal{T} = 2^X\) \, - \textbf{topologia dyskretna}
    \item \(\mathcal{T} = \{\emptyset, X\}\) \, - \textbf{topologia antydyskretna}
\end{steps}
\end{logbox}

\bigskip
\textbf{stw} \; \((X, d)\) - przestrzeń metryczna, \; \(A \subseteq X\) - otwarty \, \(\iff \, \forall_{x_0 \in A} \; \exists_{r > 0} \;\, B(x_0, r) \subseteq A\)

\bigskip
\textbf{def} \; \textbf{wnętrze} \, \(\operatorname{int} A \coloneqq \displaystyle\bigcup \{U \in \mathcal{T} : U \subseteq A\}\)

\smallskip
\textbf{def} \; \textbf{domknięcie} \, \(\overline{A} \coloneqq \displaystyle\bigcap \{F : F^c \in \mathcal{T}, \; F \supseteq A\}\)

\bigskip
\textbf{stw} \; o wnętrzu i domknięciu
\begin{podpunkty}
    \item \(\operatorname{int} A \subseteq A \subseteq \overline{A}\)
    \item \(\operatorname{int} A \in \mathcal{T}, \;\; \big(\overline{A}\big)^c \in \mathcal{T}\)
    \item \(\operatorname{int} A = \{x \in A : \exists_{U \in \mathcal{T}, \; x \in U} \;\, U \subseteq A\}\)
    \item \(\overline{A} = \big(\operatorname{int} A^c\big)^c\)
\end{podpunkty}

\smallskip
\textbf{wniosek} \; \(x \in \overline{A}^{\,c} \, \iff \, \exists_{U \in \mathcal{T}} \;\, x \in U \subseteq A^c\)

\bigskip
\textbf{def} \; \textbf{brzeg} \, \(\partial A \coloneqq \overline{A} \setminus \operatorname{int} A\)

\smallskip
\textbf{wniosek} \; \(x \in \partial A \, \iff \, \forall_{U \in \mathcal{T}, \; x \in U} \;\, U \cap A \neq \emptyset \; \land \; U \cap A^c \neq \emptyset\)

\bigskip
\textbf{tw o trichotomii} \; \((X, \mathcal{T})\) - przestrzeń topologiczna, \; \(A \subseteq X \implies X = \operatorname{int} A \sqcup \partial A \sqcup \underbrace{\overline{A}^{\,c}}_{= \, \operatorname{int} A^c}\)

\bigskip
\textbf{def} \; \(A\) jest \textbf{gęsty} w \(X\) \, \(\iff \, \overline{A} = X\)

\bigskip
\textbf{def} \; \textbf{zbiór graniczny} \; \((x_n) \subset X\), \;\; \(\operatorname{Lim}(x_n) \coloneqq \{x_0 \in X : \forall_{U \in \mathcal{T}, \; x_0 \in U} \;\, \#\{n \in \mathbb{N} : x_n \notin U\} < \infty\}\)

\bigskip
\textbf{def} \; \(x_0\) ma \textbf{przeliczalną bazę otoczeń} \, \(\iff \, \exists_{\{V_n\}_{n=1}^{\infty} \subset \mathcal{T}} \;\; \big(\forall_n \;\, x_0 \in V_n\big) \; \land \; \big(\forall_{U \in \mathcal{T}, \; x_0 \in U} \; \exists_n \;\, V_n \subseteq U\big)\)

\smallskip
\begin{logbox}
\textbf{przykład} \; w przestrzeni metrycznej \, \(V_n = B\big(x_0, \frac{1}{n}\big)\)
\end{logbox}

\bigskip
\textbf{def} \; powiemy, że przestrzeń topologiczna \((X, \mathcal{T})\) ma \textbf{przeliczalną bazę otoczeń} \, \(\iff\) każdy punkt z \(X\) ma przeliczalną bazę otoczeń

\smallskip
\textbf{wniosek} \; przestrzenie metryczne mają przeliczalne bazy otoczeń

\bigskip
\textbf{def} \; \((X, \mathcal{T})\) - \textbf{przestrzeń Hausdorffa} \, \(\iff \, \forall_{x, \hat{x} \in X, \; x \neq \hat{x}} \; \exists_{U, \hat{U} \in \mathcal{T}} \;\, x \in U, \; \hat{x} \in \hat{U}, \; U \cap \hat{U} = \emptyset\)

\smallskip
\textbf{wniosek} \; każda przestrzeń metryczna jest przestrzenią Hausdorffa

\smallskip
\textbf{fakt} \; \((X, \mathcal{T})\) - przestrzeń Hausdorffa \(\implies \forall_{(x_n) \subset X} \;\, \#\operatorname{Lim}(x_n) \leq 1\)

\bigskip
\textbf{def} \; \textbf{domknięcie ciągowe} zbioru \, \(\operatorname{scl} A \coloneqq \{x_0 \in X : \exists_{(x_n)_{n=1}^{\infty} \subset A} \;\, x_0 \in \operatorname{Lim}(x_n)\}\)

\smallskip
\textbf{tw o domknięciu ciągowym}
\begin{itemize}
    \item \(A\) - zbiór \(\implies \operatorname{scl} A \subseteq \overline{A}\)
    \item \((X, d)\) - przestrzeń metryczna, \; \(A \subseteq X \implies \operatorname{scl} A = \overline{A}\)
\end{itemize}

\smallskip
\textbf{wniosek} \; ciągowe kryterium domkniętości w przestrzeni metrycznej: \, \(A\) - domknięty \(\iff A = \operatorname{scl} A\)

\bigskip
\textbf{def} \; \((X, \mathcal{T}_X), \, (Y, \mathcal{T}_Y)\) - przestrzenie topologiczne, \; \(f: X \to Y\) - \textbf{ciągła} \, \(\iff \, \forall_{V \in \mathcal{T}_Y} \;\, f^{-1}(V) \in \mathcal{T}_X\)

\bigskip
\textbf{def} \; \(f\) spełnia \textbf{warunek Heinego} w \(x_0 \in X\) \, \(\iff \, \forall_{(x_n) \subset X} \;\, x_0 \in \operatorname{Lim}(x_n) \implies f(x_0) \in \operatorname{Lim}\big(f(x_n)\big)\) \; \((H)\)

\smallskip
\textbf{tw o równoważnych warunkach ciągłości}
\begin{itemize}
    \item \(f\) - ciągła \(\implies \forall_{x_0 \in X} \;\, f\) spełnia \((H)\) w \(x_0\)
    \item \(\mathcal{T}_X\) - metryzowalna (pochodzi od pewnej metryki), \; \(\forall_{x_0 \in X} \;\, f\) spełnia \((H)\) w \(x_0 \implies f\) - ciągła
\end{itemize}

\newpage
\textbf{def} \; \(\mathcal{B} \subseteq 2^X\), \, \textbf{topologia generowana przez bazę} \(\mathcal{B}\) \, \(\mathcal{T}(\mathcal{B}) \coloneqq \{\bigcup \mathcal{B}_0 : \mathcal{B}_0 \subseteq \mathcal{B}\}\)

\bigskip
\begin{logbox}
\textbf{przykład} \; w przestrzeni metrycznej \, \(\mathcal{T}\big(\{B(x, r) : x \in X, \; r > 0\}\big) = \mathcal{T}\)
\end{logbox}

\bigskip
\textbf{fakt} \; \(\mathcal{T}(\mathcal{B})\) - topologia \, \(\iff\)
\begin{podpunkty}
    \item \(\displaystyle\bigcup \mathcal{B} = X\)
    \item \(\forall_{U, \hat{U} \in \mathcal{B}} \;\, U \cap \hat{U} \in \mathcal{T}(\mathcal{B})\)
\end{podpunkty}

\bigskip
\begin{logbox}
\textbf{przykład} \; \(\mathcal{B} = \{[a, b] : a < b, \; a, b \in \mathbb{R}\} \implies \mathcal{T}(\mathcal{B})\) nie jest topologią, bo \([a, b] \cap [b, c] = \{b\} \notin \mathcal{T}(\mathcal{B})\)
\end{logbox}

\bigskip
\textbf{def} \; \(\mathcal{P} \subseteq 2^X\), \, \textbf{baza generowana przez podbazę} \(\mathcal{P}\) \, \(\mathcal{B}(\mathcal{P}) \coloneqq \{\bigcap_{j=1}^{n} V_j : V_1, \dots, V_n \in \mathcal{P}, \; n \in \mathbb{N}\}\)

\smallskip
\textbf{fakt} \; \(\mathcal{T}\big(\mathcal{B}(\mathcal{P})\big)\) jest zawsze topologią

\bigskip
\textbf{def} \; \((X_\alpha)_{\alpha \in A}\) - rodzina zbiorów, \; \textbf{iloczyn kartezjański} \, \(\prod \{X_\alpha : \alpha \in A\} \coloneqq \{\varphi : A \to \bigcup \{X_\alpha : \alpha \in A\} \;\; : \;\; \forall_{\alpha \in A} \;\, \varphi(\alpha) \in X_\alpha\}\)

\smallskip
\textbf{wniosek} \; \(\forall_{\alpha \in A} \;\, X_\alpha = X \implies \prod \{X_\alpha : \alpha \in A\} = X^A\)

\bigskip
\textbf{aksjomat wyboru} \; \(\forall_{\alpha \in A} \;\, X_\alpha \neq \emptyset \implies \prod \{X_\alpha : \alpha \in A\} \neq \emptyset\)


\bigskip
\bigskip
\indeks{w3-4}

\bigskip
niech \((X_\alpha, \mathcal{T}_\alpha)_{\alpha \in A}\) - przestrzenie topologiczne, \; \(X = \prod \{X_\alpha : \alpha \in A\}\)

\bigskip
\textbf{def} \; \(\alpha \in A\), \; \(U_\alpha \in \mathcal{T}_\alpha\), \; \textbf{cylinder} (otwarty) w \(X\) \, \(C_\alpha(U_\alpha) \coloneqq \{\varphi \in X : \varphi(\alpha) \in U_\alpha\}\)

\bigskip
\textbf{def} \; \textbf{topologia Tichonowa} (produktowa) na \(X\) \, \(\mathcal{T}_X \coloneqq \mathcal{T}\big(\mathcal{B}(\mathcal{P}_X)\big)\), \;\; \(\mathcal{P}_X \coloneqq \{C_\alpha(U_\alpha) : \alpha \in A, \; U_\alpha \in \mathcal{T}_\alpha\}\)

\smallskip
ozn.: \, \(\mathcal{B}_X \coloneqq \mathcal{B}(\mathcal{P}_X) = \{\bigcap_{j=1}^{n} C_{\alpha_j}(U_{\alpha_j}) : n \in \mathbb{N}_0, \; \alpha_j \in A, \; U_{\alpha_j} \in \mathcal{T}_{\alpha_j}\}\)

\smallskip
\textbf{wniosek} \; \(X = \prod_{j=1}^{n} X_j \implies \mathcal{B}_X = \{\prod_{j=1}^{n} U_j : \forall_j \;\, U_j \in \mathcal{T}_j\}\)

\bigskip
\textbf{stw} \; \((\varphi_n)_{n=0}^{\infty} \subset X \implies \Big(\varphi_0 \in \operatorname{Lim}(\varphi_n)_{n=1}^{\infty} \, \iff \, \forall_{\alpha \in A} \;\, \varphi_0(\alpha) \in \operatorname{Lim}\big(\varphi_n(\alpha)\big)_{n=1}^{\infty}\Big)\)

\bigskip
niech \((X_1, d_1), \dots, (X_{\hat{n}}, d_{\hat{n}})\) - przestrzenie metryczne, \; \(X = \prod_{j=1}^{\hat{n}} X_j\), \; \(\mathcal{T}_X\) - topologia Tichonowa

\bigskip
\textbf{def} \; \textbf{metryka produktowa} \, \(d_\Pi\big((x_j)_{j=1}^{\hat{n}}, (\hat{x}_j)_{j=1}^{\hat{n}}\big) \coloneqq \max\{d_j(x_j, \hat{x}_j) : j = 1, \dots, \hat{n}\}\)

\smallskip
\textbf{tw o równości topologii dla skończonego produktu} \; \(\mathcal{T}_X = \mathcal{T}_{d_\Pi}\)

\bigskip
\textbf{def} \; \((X, \mathcal{T})\) - przestrzeń topologiczna, \; \(Z \subseteq X\), \; \(\{U_\alpha\}_{\alpha \in A} \subset \mathcal{T}\) - \textbf{pokrycie} (otwarte) \(Z\) \, \(\iff \, Z \subseteq \bigcup \{U_\alpha : \alpha \in A\}\)

\smallskip
\textbf{def} \; \(\{U_\alpha\}_{\alpha \in \hat{A}}\) - \textbf{podpokrycie} \(\{U_\alpha\}_{\alpha \in A}\) \, \(\iff \, \hat{A} \subseteq A \; \land \; \{U_\alpha\}_{\alpha \in \hat{A}}\) - pokrycie \(Z\)

\bigskip
\textbf{def} \; \(K \subseteq X\) - \textbf{zwarty} \, \(\iff \, \forall_{\{U_\alpha\}_{\alpha \in A} \subset \mathcal{T}, \; K \subseteq \bigcup \{U_\alpha : \alpha \in A\}} \; \exists_{\hat{A} \subseteq A, \; \#\hat{A} < \infty} \;\, K \subseteq \bigcup \{U_\alpha : \alpha \in \hat{A}\}\) \, 
(dla każdego pokrycia otw. $K$ istnieje jego sk. podpokrycie)

\bigskip
\begin{logbox}
\textbf{przykład} \; \(X = \mathbb{R}\), \; \(Z = \{\frac{1}{n} : n \in \mathbb{N}\}\) \, - nie jest zwarty
\end{logbox}

\bigskip
\textbf{tw o zbiorach zwartych i domkniętych} \; \((X, \mathcal{T})\) - przestrzeń topologiczna, \; \(K \subseteq X\)
\begin{podpunkty}
    \item \(X\) - przestrzeń zwarta \; \(\land\) \; \(K\) - domknięty \(\implies K\) - zwarty
    \item \(X\) - przestrzeń Hausdorffa \; \(\land\) \; \(K\) - zwarty \(\implies K\) - domknięty
\end{podpunkty}

\bigskip
\textbf{tw Tichonowa} \; \(\forall_{\alpha \in A} \;\, (X_\alpha, \mathcal{T}_\alpha)\) - zwarta \(\implies \prod \{X_\alpha : \alpha \in A\}\) - zwarty

\bigskip
\textbf{tw o funkcjach ciągłych na zbiorach zwartych} \; \(f: X \to Y\) - ciągła, \; \(K \subseteq X\) - zwarty
\begin{podpunkty}
    \item \(f(K)\) - zwarty w \(Y\)
    \item \(Y = \mathbb{R} \implies f\) przyjmuje kresy na \(K\): \, \(\exists_{x_{\min}, x_{\max} \in K} \;\, f(x_{\min}) = \inf \{f(x) : x \in K\}, \; f(x_{\max}) = \sup \{f(x) : x \in K\}\)
    \item \((X, d_X), \, (Y, d_Y)\) - przestrzenie metryczne \(\implies f\) - \textbf{jednostajnie ciągła} na \(K\): \, \(\forall_{\varepsilon > 0} \; \exists_{\delta > 0} \; \forall_{x, \hat{x} \in K} \;\, d_X(x, \hat{x}) < \delta \implies d_Y\big(f(x), f(\hat{x})\big) < \varepsilon\)
\end{podpunkty}

\bigskip
niech \((X, d)\) - przestrzeń metryczna

\bigskip
\textbf{def} \; \(Z \subseteq X\), \; \(\varepsilon > 0\), \; \(\{x_1, \dots, x_p\} \subseteq X\) - \textbf{sieć o prześwicie} \(\varepsilon\) dla \(Z\) \, \(\iff \, Z \subseteq \bigcup_{j=1}^{p} B(x_j, \varepsilon)\)

\smallskip
\textbf{def} \; \(Z\) - \textbf{całkowicie ograniczony} \, \(\iff \, \forall_{\varepsilon > 0} \; \exists\) sieć o prześwicie \(\varepsilon\) dla \(Z\)

\bigskip
\begin{logbox}
\textbf{przykład} \; \(\#X = \infty\), \; \(d(x, y) = 1 - \delta_x^y\) \,(metryka dyskretna) \, - \(X\) ograniczony, ale nie całkowicie ograniczony
\end{logbox}

\bigskip
\textbf{fakt} \; \(\exists\) sieć o prześwicie \(\varepsilon\) dla \(Z \implies \exists\) sieć \(\{x_1, \dots, x_p\} \subseteq Z\) o prześwicie \(2\varepsilon\) dla \(Z\)

\bigskip
\textbf{def} \; \((x_n)\) - \textbf{zdystansowany} o \(\varepsilon > 0\) \, \(\iff \, \forall_{n, \hat{n}, \; n \neq \hat{n}} \;\, d(x_n, x_{\hat{n}}) \geq \varepsilon\)

\smallskip
\textbf{def} \; \((x_n)\) - \textbf{zdystansowany} \, \(\iff \, \exists_{\varepsilon > 0} \;\, (x_n)\) - zdystansowany o \(\varepsilon\)

\bigskip
\textbf{kryterium całkowitej ograniczoności} \; \(Z\) - całkowicie ograniczony \(\iff\) nie istnieje ciąg zdystansowany \((x_n) \subset Z\)

\newpage
\indeks{w5}

\bigskip
\textbf{def} \; \((X, \mathcal{T})\) - przestrzeń topologiczna, \; \((x_n) \subset X\), \; \(\operatorname{Acc}(x_n)_{n=1}^{\infty} \coloneqq \{x \in X : \forall_{U \in \mathcal{T}, \; x \in U} \;\, \#\{n \in \mathbb{N} : x_n \in U\} = \infty\}\)

\smallskip
\textbf{def} \; \(x_0 \in X\) - \textbf{punkt skupienia} ciągu \((x_n)\) \, \(\iff \, x_0 \in \operatorname{Acc}(x_n)\)

\smallskip
\textbf{fakt} \; \(\operatorname{Lim}(x_n) \subseteq \operatorname{Acc}(x_n)\)

\bigskip
\begin{logbox}
\textbf{przykład} \; \(x_n = (-1)^n\) \, - \(\operatorname{Lim}(x_n) = \emptyset \subseteq \{-1, 1\} = \operatorname{Acc}(x_n)\)
\end{logbox}

\bigskip
\textbf{tw o zwartości w przestrzeniach metrycznych} \; \((X, d)\) - przestrzeń metryczna, \; \(K \subseteq X\), \; NWSR:
\begin{podpunkty}
    \item \(K\) - zwarty
    \item \(\forall_{(x_n) \subset K} \;\, \operatorname{Acc}(x_n) \cap K \neq \emptyset\) \, (każdy ciąg o wyrazach z $K$ ma punkt skupienia należący do $K$)
    \item \(\forall_{(x_n) \subset K} \; \exists_{(x_{n_k}) \text{ - podciąg}} \;\, \operatorname{Lim}(x_{n_k}) \cap K \neq \emptyset\) \, (każdy ciąg o wyrazach z $K$ ma podciąg zbieżny do granicy w $K$)
    \item \(K\) - całkowicie ograniczony \(\land \; (K, d)\) - zupełna
\end{podpunkty}

\bigskip
\textbf{tw o całkowitej ograniczoności w przestrzeniach} \(\mathbb{R}^n\) \; \(K \subseteq \mathbb{R}^n\) - ograniczony \(\implies K\) - całkowicie ograniczony

\bigskip
\textbf{tw Heinego-Borela} \; \(K \subseteq \mathbb{R}^n\) - zwarty \, \(\iff \, K\) - domknięty \(\land \; K\) - ograniczony


\bigskip
\bigskip
\indeks{w6-7}

\bigskip
\temat{teoria miary}


\end{document}
